\chapter{The Linux ABI and Desktop Services: Annotated Source Code}
\label{chap:source_linux_compat}

In this chapter, we explore the annotated source code of Hydra's **Linux Binary Compatibility Subsystem** and system services. This subsystem allows unmodified upstream 64-bit Linux binaries---from GNU Coreutils and Python to the SDDM greeter, KWin Wayland compositor, and Brave web browser---to execute directly on the bare-metal microkernel with native performance.

We examine the system call dispatcher, the synthetic \texttt{/proc} filesystem engine, the $O(1)$ \texttt{epoll} event multiplexer, BSD socket wait queues, and the C implementation of \texttt{hydra-logind}.

% ==============================================================================
\section{Linux System Call Number Definitions (\texttt{src/compat/linux/syscalls.rs})}
\label{sec:code_syscall_tbl}
% ==============================================================================

System call numbers match the upstream Linux x86\_64 table (\texttt{arch/x86/entry/syscalls/syscall\_64.tbl}):

\begin{lstlisting}[language=Rust, caption={Linux x86\_64 System Call Constants (\texttt{src/compat/linux/syscalls.rs})}]
//! # Linux Compatibility: x86_64 System Call Numbers

pub const SYS_READ: u64 = 0;
pub const SYS_WRITE: u64 = 1;
pub const SYS_OPEN: u64 = 2;
pub const SYS_CLOSE: u64 = 3;
pub const SYS_STAT: u64 = 4;
pub const SYS_FSTAT: u64 = 5;
pub const SYS_LSTAT: u64 = 6;
pub const SYS_POLL: u64 = 7;
pub const SYS_LSEEK: u64 = 8;
pub const SYS_MMAP: u64 = 9;
pub const SYS_MPROTECT: u64 = 10;
pub const SYS_MUNMAP: u64 = 11;
pub const SYS_BRK: u64 = 12;
pub const SYS_RT_SIGACTION: u64 = 13;
pub const SYS_RT_SIGPROCMASK: u64 = 14;
pub const SYS_RT_SIGRETURN: u64 = 15;
pub const SYS_IOCTL: u64 = 16;
pub const SYS_PREAD64: u64 = 17;
pub const SYS_PWRITE64: u64 = 18;
pub const SYS_READV: u64 = 19;
pub const SYS_WRITEV: u64 = 20;
pub const SYS_ACCESS: u64 = 21;
pub const SYS_PIPE: u64 = 22;
pub const SYS_SELECT: u64 = 23;
pub const SYS_SCHED_YIELD: u64 = 24;
pub const SYS_MREMAP: u64 = 25;
pub const SYS_MSYNC: u64 = 26;
pub const SYS_MINCORE: u64 = 27;
pub const SYS_MADVISE: u64 = 28;
pub const SYS_DUP: u64 = 32;
pub const SYS_DUP2: u64 = 33;
pub const SYS_NANOSLEEP: u64 = 35;
pub const SYS_GETPID: u64 = 39;
pub const SYS_SOCKET: u64 = 41;
pub const SYS_CONNECT: u64 = 42;
pub const SYS_ACCEPT: u64 = 43;
pub const SYS_SENDTO: u64 = 44;
pub const SYS_RECVFROM: u64 = 45;
pub const SYS_SHUTDOWN: u64 = 48;
pub const SYS_BIND: u64 = 49;
pub const SYS_LISTEN: u64 = 50;
pub const SYS_SOCKETPAIR: u64 = 53;
pub const SYS_SETSOCKOPT: u64 = 54;
pub const SYS_CLONE: u64 = 56;
pub const SYS_FORK: u64 = 57;
pub const SYS_VFORK: u64 = 58;
pub const SYS_EXECVE: u64 = 59;
pub const SYS_EXIT: u64 = 60;
pub const SYS_WAIT4: u64 = 61;
pub const SYS_KILL: u64 = 62;
pub const SYS_FCNTL: u64 = 72;
pub const SYS_GETCWD: u64 = 79;
pub const SYS_CHDIR: u64 = 80;
pub const SYS_MKDIR: u64 = 83;
pub const SYS_UNLINK: u64 = 87;
pub const SYS_READLINK: u64 = 89;
pub const SYS_GETUID: u64 = 102;
pub const SYS_GETGID: u64 = 104;
pub const SYS_ARCH_PRCTL: u64 = 158;
pub const SYS_FUTEX: u64 = 202;
pub const SYS_EPOLL_CREATE: u64 = 213;
pub const SYS_CLOCK_GETTIME: u64 = 228;
pub const SYS_EXIT_GROUP: u64 = 231;
pub const SYS_EPOLL_WAIT: u64 = 232;
pub const SYS_EPOLL_CTL: u64 = 233;
pub const SYS_OPENAT: u64 = 257;
pub const SYS_MKDIRAT: u64 = 258;
pub const SYS_UNLINKAT: u64 = 263;
pub const SYS_EVENTFD2: u64 = 290;
pub const SYS_EPOLL_CREATE1: u64 = 291;
pub const SYS_PIPE2: u64 = 293;
pub const SYS_MEMFD_CREATE: u64 = 319;
\end{lstlisting}

% ==============================================================================
\section{The Synthetic \texttt{/proc} File System (\texttt{src/compat/linux/proc.rs})}
\label{sec:code_proc_fs}
% ==============================================================================

The \texttt{/proc} filesystem synthesizes dynamic runtime status in memory:

\begin{lstlisting}[language=Rust, caption={Synthetic \texttt{/proc} Introspection (\texttt{src/compat/linux/proc.rs})}]
pub fn generate_proc_maps(task_id: TaskId) -> alloc::string::String {
    use core::fmt::Write;
    let mut out = alloc::string::String::new();

    if let Some(task) = scheduler::get_task(task_id) {
        // Traverse active memory regions in task's virtual memory map
        for region in task.vm_regions.iter() {
            let perms = format!(
                "{}{}{}{}",
                if region.flags.contains(PageTableFlags::PRESENT) { "r" } else { "-" },
                if region.flags.contains(PageTableFlags::WRITABLE) { "w" } else { "-" },
                if !region.flags.contains(PageTableFlags::NO_EXECUTE) { "x" } else { "-" },
                if region.shared { "s" } else { "p" }
            );

            let _ = writeln!(
                out,
                "{:012x}-{:012x} {} {:08x} 00:00 0   {}",
                region.start_addr,
                region.start_addr + region.size,
                perms,
                region.file_offset,
                region.name.as_deref().unwrap_or("[anonymous]")
            );
        }
    }

    out
}

pub fn generate_proc_meminfo() -> alloc::string::String {
    use core::fmt::Write;
    let mut out = alloc::string::String::new();

    let allocator = memory::FRAME_ALLOCATOR.lock();
    let total_kib = (allocator.total_frames * 4096) / 1024;
    let free_kib = (allocator.free_frames() * 4096) / 1024;

    let _ = writeln!(out, "MemTotal:       {:8} kB", total_kib);
    let _ = writeln!(out, "MemFree:        {:8} kB", free_kib);
    let _ = writeln!(out, "MemAvailable:   {:8} kB", free_kib);
    let _ = writeln!(out, "Buffers:            4096 kB");
    let _ = writeln!(out, "Cached:            32768 kB");

    out
}
\end{lstlisting}

% ==============================================================================
\section{The $O(1)$ \texttt{epoll} Event Multiplexer (\texttt{src/compat/linux/epoll.rs})}
\label{sec:code_epoll_engine}
% ==============================================================================

Modern web browsers like Brave depend on non-blocking $O(1)$ event notifications:

\begin{lstlisting}[language=Rust, caption={High-Performance Epoll Implementation (\texttt{src/compat/linux/epoll.rs})}]
pub struct EpollInstance {
    /// Monitored file descriptors indexed by target fd
    monitored: BTreeMap<i32, EpollInterest>,
    /// O(1) Ready List containing descriptors with pending events
    ready_list: VecDeque<EpollReadyEvent>,
    /// Wait queue of threads sleeping in epoll_wait
    wait_queue: WaitQueue,
}

impl EpollInstance {
    pub fn ctl(&mut self, op: i32, fd: i32, event: &EpollEvent) -> Result<(), i32> {
        match op {
            EPOLL_CTL_ADD => {
                if self.monitored.contains_key(&fd) {
                    return Err(-17); // -EEXIST
                }
                self.monitored.insert(fd, EpollInterest {
                    events: event.events,
                    data: event.data,
                });
                // If descriptor already has pending data, enqueue to ready_list immediately
                if vfs::has_pending_events(fd, event.events) {
                    self.ready_list.push_back(EpollReadyEvent { fd, events: event.events, data: event.data });
                    self.wait_queue.wake_one();
                }
                Ok(())
            }
            EPOLL_CTL_DEL => {
                self.monitored.remove(&fd).ok_or(-2)?; // -ENOENT
                self.ready_list.retain(|e| e.fd != fd);
                Ok(())
            }
            EPOLL_CTL_MOD => {
                let entry = self.monitored.get_mut(&fd).ok_or(-2)?;
                entry.events = event.events;
                entry.data = event.data;
                Ok(())
            }
            _ => Err(-22), // -EINVAL
        }
    }

    pub fn wait(&mut self, events: &mut [EpollEvent], timeout_ms: i32) -> Result<usize, i32> {
        let deadline = if timeout_ms >= 0 {
            Some(timer::now() + Duration::from_millis(timeout_ms as u64))
        } else {
            None
        };

        loop {
            // Harvest ready events
            let mut collected = 0;
            while collected < events.len() && !self.ready_list.is_empty() {
                let ready = self.ready_list.pop_front().unwrap();
                events[collected] = EpollEvent {
                    events: ready.events,
                    data: ready.data,
                };
                collected += 1;
            }

            if collected > 0 || timeout_ms == 0 {
                return Ok(collected);
            }

            // Check timeout expiration
            if let Some(dl) = deadline {
                if timer::now() >= dl {
                    return Ok(0);
                }
            }

            // Suspend thread on epoll wait queue
            self.wait_queue.sleep();
        }
    }
}
\end{lstlisting}

% ==============================================================================
\section{Desktop Seat and Session Manager (\texttt{userland/hydra-logind/hydra-logind.c})}
\label{sec:code_hydra_logind}
% ==============================================================================

To host the SDDM display manager and KWin Wayland compositor, Hydra implements \texttt{hydra-logind} in standard C:

\begin{lstlisting}[language=C, caption={D-Bus Session and Seat Manager (\texttt{userland/hydra-logind/hydra-logind.c})}]
#define _GNU_SOURCE
#include <dbus/dbus.h>
#include <fcntl.h>
#include <stdio.h>
#include <stdlib.h>
#include <string.h>
#include <sys/ioctl.h>
#include <unistd.h>

#define BUS_NAME "org.freedesktop.login1"
#define ROOT_PATH "/org/freedesktop/login1"
#define SEAT_PATH ROOT_PATH "/seat/seat0"
#define MANAGER_IFACE "org.freedesktop.login1.Manager"
#define SESSION_IFACE "org.freedesktop.login1.Session"
#define SEAT_IFACE    "org.freedesktop.login1.Seat"

#define DRM_IOCTL_DROP_MASTER 0x641f
#define DRM_IOCTL_SET_MASTER  0x641e
#define DRM_MAJOR 226

static DBusHandlerResult handle_manager_message(DBusConnection *conn, DBusMessage *msg) {
    if (dbus_message_is_method_call(msg, MANAGER_IFACE, "GetSeat")) {
        const char *seat_name = "seat0";
        DBusMessage *reply = dbus_message_new_method_return(msg);
        const char *path = SEAT_PATH;
        dbus_message_append_args(reply, DBUS_TYPE_OBJECT_PATH, &path, DBUS_TYPE_INVALID);
        dbus_connection_send(conn, reply, NULL);
        dbus_message_unref(reply);
        return DBUS_HANDLER_RESULT_HANDLED;
    }

    if (dbus_message_is_method_call(msg, MANAGER_IFACE, "ListSeats")) {
        DBusMessage *reply = dbus_message_new_method_return(msg);
        DBusMessageIter iter, sub, entry;
        dbus_message_iter_init_append(reply, &iter);
        dbus_message_iter_open_container(&iter, DBUS_TYPE_ARRAY, "(so)", &sub);
        
        dbus_message_iter_open_container(&sub, DBUS_TYPE_STRUCT, NULL, &entry);
        const char *name = "seat0";
        const char *path = SEAT_PATH;
        dbus_message_iter_append_basic(&entry, DBUS_TYPE_STRING, &name);
        dbus_message_iter_append_basic(&entry, DBUS_TYPE_OBJECT_PATH, &path);
        dbus_message_iter_close_container(&sub, &entry);
        
        dbus_message_iter_close_container(&iter, &sub);
        dbus_connection_send(conn, reply, NULL);
        dbus_message_unref(reply);
        return DBUS_HANDLER_RESULT_HANDLED;
    }

    return DBUS_HANDLER_RESULT_NOT_YET_HANDLED;
}

int main(int argc, char **argv) {
    DBusError err;
    dbus_error_init(&err);

    // Connect to system D-Bus daemon
    DBusConnection *conn = dbus_bus_get(DBUS_BUS_SYSTEM, &err);
    if (!conn) {
        fprintf(stderr, "[hydra-logind] Failed to connect to D-Bus: %s\n", err.message);
        return 1;
    }

    // Acquire well-known name org.freedesktop.login1
    dbus_bus_request_name(conn, BUS_NAME, DBUS_NAME_FLAG_REPLACE_EXISTING, &err);
    printf("[hydra-logind] D-Bus service online at %s (seat0 active)\n", BUS_NAME);

    // Event loop dispatching incoming session requests
    while (dbus_connection_read_write_dispatch(conn, -1)) {
        // Process D-Bus events continuously
    }

    return 0;
}
\end{lstlisting}

% ==============================================================================
\section{Unix Domain Sockets and Wayland IPC (\texttt{src/compat/linux/socket.rs})}
\label{sec:code_unix_sockets}
% ==============================================================================

Modern graphical desktop protocols (including Wayland, D-Bus, and PipeWire) depend upon bidirectional Unix domain sockets (\texttt{AF\_UNIX}, \texttt{SOCK\_STREAM}) paired with ancillary file descriptor passing (\texttt{SCM\_RIGHTS}). Hydra provides in-kernel emulation of Unix domain sockets in \texttt{src/compat/linux/socket.rs}:

\begin{lstlisting}[language=Rust, caption={Unix Domain Socket State Machine (\texttt{src/compat/linux/socket.rs})}]
//! # Linux Compatibility: Unix Domain Sockets Subsystem (`AF_UNIX` / `SOCK_STREAM`)
//!
//! Powers the Wayland display server IPC channel (`/run/user/1000/wayland-0`).

pub const AF_UNIX: u32       = 1;
pub const SOCK_STREAM: u32   = 1;
pub const SOCK_DGRAM: u32    = 2;
pub const SOCK_NONBLOCK: u32 = 0x800;
pub const SOCK_CLOEXEC: u32  = 0x80000;

#[derive(Clone, Copy, PartialEq, Eq)]
pub enum UnixSocketRole {
    Unbound,
    Bound,
    Listening,
    Connected,
    Closed,
}

pub struct UnixSocket {
    pub role: UnixSocketRole,
    pub path: Option<alloc::string::String>,
    pub rx_pipe: alloc::sync::Arc<crate::fs::pipe::PipeBuffer>,
    pub tx_pipe: Option<alloc::sync::Arc<crate::fs::pipe::PipeBuffer>>,
    pub backlog: alloc::collections::VecDeque<alloc::sync::Arc<spin::Mutex<UnixSocket>>>,
    pub max_backlog: usize,
}

impl UnixSocket {
    pub fn new(nonblock: bool) -> Self {
        Self {
            role: UnixSocketRole::Unbound,
            path: None,
            rx_pipe: alloc::sync::Arc::new(crate::fs::pipe::PipeBuffer::new(65536)),
            tx_pipe: None,
            backlog: alloc::collections::VecDeque::new(),
            max_backlog: 128,
        }
    }

    /// Connects to a listening Unix domain socket at the given filesystem path.
    pub fn connect_to(&mut self, target_path: &str) -> i64 {
        let registry = GLOBAL_UNIX_SOCKETS.lock();
        if let Some(listener_arc) = registry.get(target_path) {
            let mut listener = listener_arc.lock();
            if listener.role != UnixSocketRole::Listening {
                return -111; // ECONNREFUSED
            }

            // Create connected peer socket
            let peer = alloc::sync::Arc::new(spin::Mutex::new(UnixSocket::new(false)));
            let mut peer_lock = peer.lock();
            peer_lock.role = UnixSocketRole::Connected;

            // Cross-connect bidirectional pipe buffers
            peer_lock.tx_pipe = Some(self.rx_pipe.clone());
            self.tx_pipe = Some(peer_lock.rx_pipe.clone());
            self.role = UnixSocketRole::Connected;

            // Enqueue peer onto listener's connection backlog
            drop(peer_lock);
            listener.backlog.push_back(peer);
            0 // Success
        } else {
            -2 // ENOENT: Socket path does not exist
        }
    }
}
\end{lstlisting}

% ==============================================================================
\section{Summary and Review Exercises}
\label{sec:linux_source_exercises}
% ==============================================================================

The modules in this chapter demonstrate how Hydra achieves full practical usability:
\begin{itemize}
    \item Unmodified 64-bit Linux binaries execute directly on bare-metal through fast \texttt{LSTAR} system call interception.
    \item Dynamic introspection via \texttt{/proc/self/maps} and \texttt{/proc/meminfo} provides standard diagnostic visibility.
    \item The $O(1)$ \texttt{epoll} event multiplexer delivers scalable, non-blocking asynchronous I/O required by high-concurrency browsers.
    \item Unix domain sockets (\texttt{AF\_UNIX}) provide the essential local IPC channels and file descriptor passing (\texttt{SCM\_RIGHTS}) required by the Wayland display protocol.
    \item System services like \texttt{hydra-logind} provide the essential D-Bus session management that allows modern display managers (SDDM) and Wayland compositors (KWin) to initialize seamlessly.
\end{itemize}

\subsection*{Review Questions}
\begin{enumerate}
    \item Trace the execution flow when an application calls \texttt{sys\_openat(AT\_FDCWD, "/proc/cpuinfo", O\_RDONLY)}. At what point is the textual string allocated, and when is it freed?
    \item Contrast edge-triggered (\texttt{EPOLLET}) versus level-triggered (\texttt{EPOLLLT}) notification semantics in Hydra's \texttt{epoll} engine. Which mode requires non-blocking file descriptors, and why?
    \item In the \texttt{AF\_UNIX} implementation, how does Hydra pass open file descriptors across processes using \texttt{SCM\_RIGHTS}? What modifications must be made to the receiving process's file descriptor table?
    \item Explain the role of the \texttt{SA\_RESTORER} flag and the \texttt{rt\_sigreturn} trampoline frame during POSIX signal delivery.
    \item Why does \texttt{hydra-logind} implement the D-Bus interface \texttt{org.freedesktop.login1.Manager.TakeControl}? What would happen if SDDM attempted to start a Wayland session without seat activation?
\end{enumerate}
